% !TEX root =  master.tex

\newcommand{\weigert}{Julian Weigert\xspace}
% -------------------------------------------------------
% umfrageteilnehmende Person
%
% ------------------------------------------------------- 
\newpage
\cvsect{\weigert} %sich zu arbeiten.
\begin{minipage}[t]{0.5\textwidth}
	\vspace{-3.6cm}
	\renewcommand{\arraystretch}{1.5}
	\begin{tabular}{l l}
		Name: & \weigert \\
		Alter: & 24 \\
		Tätigkeit: & Student der BWL \\
		Skills: & 2/5 \hspace{-1cm} \begin{barchart}{5.0}
			\baritemNL{}{2}
		\end{barchart} \\
	\end{tabular}
\end{minipage}
\hfill
\begin{minipage}[t]{0.4\textwidth}
	\flushright
	\includegraphics[width=0.70\textwidth]{img/personas/julian}
\end{minipage}
\autocite{rf-unsplash-student}

Das möchte ich gerne haben:
\begin{itemize}
	\item Umfragen anonym durchführen
    \item Umfragen auf seinem Smartphone durchführen
    \item Umfragen auch später durchführen
\end{itemize}

Mit der dritten Persona, \weigert, wird ein technisch weniger begabter Student einer Hochschule beschrieben.
\weigert nimmt primär die Rolle des Umfrageteilnehmers ein, wobei er die Umfragen gerne an seinem mobilen Endgeräten bearbeiten möchte.
Durch seinen eng getakteten Zeitplan möchte er Umfragen auch Tage nach der Herausgabe des Links beantworten können.
